\chapter{The Machine}
% OUTLINE Ch4 "The Machine (tensors and forms)". Laws: CITATION (MTW 1973
% slots pedagogy; Cauchy 1822 stress; Cartan 1899 forms + the general
% Stokes; the classical Stokes lineage told with attribution honesty:
% Green 1828, Gauss/Ostrogradsky divergence, Kelvin->Stokes 1850/1854 exam,
% trade's name stuck to the examiner), LENS (zero device appearances),
% Vol-1 laws (step-not-move watched). THE NS SEED-AS-DEBT: one fenced
% sentence at the Cauchy paragraph, per the outline. No metric anywhere:
% the exterior derivative's metric-freedom is the no-ruler theme's beat.
% Gate: Kodo per-claim citation.

The last chapter ended with a promise of multiplication, and the cleanest
way to keep it is with a picture the physicist may already own. In the
great textbook of gravitation that Misner, Thorne, and Wheeler published in
1973, a tensor is introduced as a machine --- a device with input slots.
Into some slots you insert arrows; into others, read-outs; the machine
returns a single number, and it is linear in every slot separately: double
one input and the output doubles, sum two inputs in a slot and the outputs
sum. That is the entire definition, and this chapter asks the reader to
take it literally, because everything else follows from the slots. A
covector was already such a machine, with one arrow-slot. A tensor is the
same idea grown up: a machine of type $(k,l)$ takes $l$ arrows and $k$
read-outs, and answers, multilinearly, with a number.

The bookkeeping comes for free, and the reader should notice that nothing
new is decreed. Feed the machine basis elements of a chart and record the
answers: those are the components, one number per choice of inputs, the
many-indexed arrays the physicist manipulates daily. Change observers, and
each slot's inputs transform by the law the last chapter derived for their
species --- arrows by the Jacobian, read-outs by its inverse --- so the
machine's components transform slot by slot, one factor of the dictionary
per index. The general tensor transformation law, the one that fills a
blackboard when written in full, is not a new law at all; it is Chapter 3's
two laws, composed. The calculus of such component-systems is again Ricci
and Levi-Civita's absolute differential calculus of 1901; the machine
picture and the component picture are two presentations of one object, and the
book will use whichever is honest and shorter.

The oldest tensor in physics was not proposed by a geometer; it was found
inside a loaded body. Cauchy, in 1822, analyzing the internal forces of
continuous matter, established the fact that deserves restating in the
machine's language: cut a small surface element anywhere inside a stressed
body, and the traction across it --- the force per area the material on one
side exerts on the other --- depends \emph{linearly} on the element's
orientation. Insert the surface's normal direction into a slot; receive the
traction. That machine is the stress tensor, the object from which the
whole subject took its name --- \emph{tensor}, from tension --- and it is
the reason the word belongs to physics by birthright rather than by loan.

\begin{fence}
A debt acknowledged in passing, not paid. The stress tensor is the native
object of the fluid description --- the fourth station of this book's
preface, the sounding promised in the founding text of the apparatus and
not yet taken. It appears here as the historical headwater of the tensor
concept, and this chapter spends nothing from that owed account.
\end{fence}

A tensor field is then a tensor at every point --- a machine stationed at
each event of the manifold, its components varying smoothly in every chart
--- and with that phrase the reader already holds the form in which physics
states its laws: the fields of the trade are tensor fields, stresses and
fluxes and, chapters from now, curvatures.

Among the machines, one family earns a chapter-within-a-chapter: the
machines that care about orientation. Ask a machine with several
arrow-slots to change sign whenever two of its arrow inputs are exchanged,
and what it measures is no longer a list of directions but an oriented
patch --- the parallelogram two arrows span, the cell three arrows span,
sign flipping when the orientation flips. These alternating machines are
the differential forms, and their calculus --- the exterior calculus ---
was built by \'Elie Cartan beginning in 1899. Two facts about it carry this
book's longest theme another step. First, forms are the natural integrands:
what one integrates over a curve, a surface, a region, is precisely a form
of matching degree, because integration is nothing but summing a machine's
readings over the oriented cells that tile the domain. Second --- and the
reader should hear the theme --- Cartan's derivative of a form, the
exterior derivative, is defined with no metric and no further apparatus at
all: from the treaty's smooth structure alone, each $k$-form has a
$(k{+}1)$-form derivative, measuring how the machine's readings on a cell's
boundary fail to cancel. No ruler was admitted, and a full differentiation
arrived anyway. The inventory of what stands without lengths now includes
an integral and a derivative.

Those two constructions meet in the one theorem this chapter states in
full, because the reader has known its special cases since their second
year: the integral of a form's exterior derivative over a region equals the
integral of the form over the region's boundary. The boundary reads the
interior. Every classical integral identity is this sentence in a
particular costume, and the costumes have a history worth telling with the
credit put right, since this series has a habit about that. The plane case
is Green's, published in 1828 in a privately printed essay. The divergence
case belongs to Gauss and to Ostrogradsky, who gave the general form in
1831. The curl case was communicated by Kelvin in a letter to Stokes in
1850, and Stokes set it as a question in the Smith's Prize examination of
1854 --- and the trade, reading the exam rather than the letter, attached
the examiner's name to the examinee's theorem, where it has stuck ever
since. The general statement, uniting all of them for forms of every
degree, belongs to the exterior calculus and is Cartan's. This book calls
the general theorem Stokes' theorem, as the trade does, and records here
whose it is.

Take the inventory once more, because it has grown quietly and the reader
should audit it. Machines at a point; fields of machines; among them the
oriented family, with an integral calculus and a derivative of their own;
and one master identity binding derivative to boundary. Every construction
cited, none original, and --- the standing astonishment --- still no
metric: nothing yet measures a length, and none of the above required one.

What the machines cannot yet do is talk to their neighbors. A tensor at
this point and a tensor at the next point live in different rooms ---
different tangent spaces --- and nothing constructed so far compares them:
subtraction of neighbors, the first gesture of differentiation for arrows
and general machines, is not yet defined. The missing structure is a rule
for carrying machines from room to room, it is genuinely extra --- no
treaty clause supplies it --- and choosing it honestly is the next
chapter's whole business.
